\begin{bt}%[1D9H1-2]
% nguon: Archive/2024/TOAN_11_XE/TOAN_11_XE/data4/11C8/KNTT-8.1.tex #1226

	Hộp thứ nhất chứa $3$ tấm thẻ cùng loại được đánh số lần lượt từ $1$ đến $3$. Hộp thứ hai chứa $5$ tấm thẻ cùng loại được đánh số lần lượt từ $1$ đến $5$. Lấy ra ngẫu nhiên từ mỗi hộp $1$ thẻ. Gọi $A$ là biến cố \lq\lq Tổng các số ghi trên $2$ thẻ bằng $6$\rq\rq, $B$ là biến cố \lq\lq Tích các số ghi trên $2$ thẻ là số lẻ\rq\rq. 
	\begin{enumerate}
		\item [a)] Hãy viết tập hợp mô tả biến cố $A B$ và tính $P(A B)$.
		\item [b)] Hãy tìm một biến cố khác rỗng và xung khắc với cả hai biến cố $A$ và $B$.	
	\end{enumerate}
	\loigiai{
		\begin{enumerate}
			\item [a)] Hãy viết tập hợp mô tả biến cố $A B$ và tính $P(A B)$.\\
			Hộp thứ nhất chứa $3$ tấm thẻ cùng loại được đánh số $1; 2; 3$.\\
			Hộp thứ hai chứa $5$ tấm thẻ cùng loại được đánh số lần lượt từ $1; 2; 3; 4; 5$.\\
			Số phần tử không gian mẫu là $3\cdot 5=15$.\\
			Gọi $A$ là biến cố \lq\lq Tổng các số ghi trên $2$ thẻ bằng $6$\rq\rq $\,$nên $A=\{(1;5);(2;4);(3;3)\}$. $P(A)=\dfrac{3}{15}$.\\
			Gọi $B$ là biến cố \lq\lq Tích các số ghi trên $2$ thẻ là số lẻ\rq\rq $\,$nên $B=\{(1;1);(1;3);(1;5);(3;1);(3;3);(3;5)\}$. $P(B)=\dfrac{6}{15}=\dfrac{2}{5}$. \\
			Vậy $P(A B)=P(A)P(B)=\dfrac{1}{5}\cdot \dfrac{2}{5}=\dfrac{2}{25}$.
			\item [b)] Một biến cố khác rỗng và xung khắc với cả hai biến cố $A$ và $B$ là \\
			$C$ là biến cố \lq\lq Tích các số ghi trên $2$ thẻ là số chẵn và tổng khác 6\rq\rq .\\
			Suy ra $C=\{(1;2);(1;4);(2;1);(2;2);(2;3);(2;5);(3;2);(3;4)\}$.
		\end{enumerate}
	}
\end{bt}

\begin{bt}%[1D9H2-2]
% nguon: Archive/2024/TOAN_11_XE/TOAN_11_XE/data4/11C8/KNTT-8.1.tex #1227

	Một hộp chứa $21$ tấm thẻ cùng loại được đánh số từ $1$ đến $21$. Chọn ra ngẫu nhiên $1$ thẻ từ hộp. Gọi $A$ là biến cố \lq\lq Số ghi trên thẻ được chọn chia hết cho $2$\rq\rq, $B$ là biến cố \lq\lq Số ghi trên thẻ được chọn chia hết cho $3$\rq\rq.
	\begin{enumerate}
		\item [a)] Hãy mô tả bằng lời biến cố $A B$.
		\item [b)] Hai biến cố $A$ và $B$ có độc lập không? Tại sao?
	\end{enumerate}
	\loigiai{
		\begin{enumerate}
			\item [a)] $A$ là biến cố \lq\lq Số ghi trên thẻ được chọn chia hết cho $2$\rq\rq $\,$nên $A=\{2;4;6;8;10;12;14;16;18;20\}$.\\
			$B$ là biến cố \lq\lq Số ghi trên thẻ được chọn chia hết cho $3$\rq\rq $\,$nên $B=\{3;6;9;12;15;18;21\}$.\\
			Suy ra biến cố $AB=\{6;12;18\}$. \\
			Vậy biến cố $AB$ là \lq\lq Số ghi trên thẻ được chọn chia hết cho 6\rq\rq.
			\item [b)] 
			Xác suất $P(A)=\dfrac{10}{21}$;	Xác suất $P(B)=\dfrac{7}{21}=\dfrac{1}{3}$; Xác suất $P(AB)=\dfrac{3}{21}=\dfrac{1}{7}$.\\
			Ta có $P(AB)=\dfrac{1}{7}\neq \dfrac{10}{21}\cdot \dfrac{1}{3}=\dfrac{10}{63}=P(A)P(B)$ nên hai biến cố $A$ và $B$ không độc lập.
		\end{enumerate}	
	}
\end{bt}
